% \begin{abstract}
%     Large language models (LLMs) typically allocate significant attention weights to the initial tokens even though they are not semantically important, which is called \textbf{attention sink}. This phenomenon is well-adopted in the scenarios of streaming generation, long context understanding, KV cache optimization, inference acceleration, etc. Despite its widespread applications, a deep understanding of attention sink in LLMs is still absent. Looking into this, we observe that attention sink already appears in the pre-training stage, which motivates us to investigate how optimization, data, auto-regressive objectives, and model parameterizations affect attention sink in LLMs. Besides a comprehensive study in quantifying the effects of the above factors, we highlight that attention sink appears after LLMs are trained on enough diverse data with effective optimization. The reason why attention sink appears in the first token is attributed to that only the first token could be ``looked at'' by all other tokens.  Meanwhile, we show that LLMs with softmax attention need key bias terms in the computation of attention even though the corresponding values are zeros, resulting in the attention sink. This motivates us to replace softmax attention with other attention operations, e.g., sigmoid attention, which results in LLMs without attention sink. 
%     % To fill in the gap, we conduct a comprehensive study to disclose how various factors affect attention sink, with a focus on LLMs pre-training.  
% \end{abstract}

\begin{abstract}
    Auto-regressive Language Models (LMs) assign significant attention to the first token, even if it is not semantically important, which is known as \textbf{attention sink}. This phenomenon has been widely adopted in applications such as streaming/long context generation, KV cache optimization, inference acceleration, model quantization, and others. Despite its widespread use, a deep understanding of attention sink in LMs is still lacking. In this work, we first demonstrate that attention sinks exist universally in auto-regressive LMs with various inputs, even in small models. Furthermore, attention sink is observed to emerge during the LM pre-training, motivating us to investigate how {\color{cc1}{optimization}}, {\color{cc2}{data distribution}}, {\color{cc3}{loss function}}, and {\color{cc4}{model architecture}} in LM pre-training influence its emergence. We highlight that attention sink emerges after effective optimization on sufficient training data. The sink position is highly correlated with the loss function and data distribution. Most importantly, we find that attention sink acts more like key biases, \emph{storing extra attention scores}, which could be non-informative and not contribute to the value computation. We also observe that this phenomenon (at least partially) stems from tokens' inner dependence on attention scores as a result of softmax normalization. After relaxing such dependence by replacing softmax attention with other attention operations, such as sigmoid attention without normalization, attention sinks do not emerge in LMs up to 1B parameters. The code is available at \href{https://github.com/sail-sg/Attention-Sink}{https://github.com/sail-sg/Attention-Sink}.\looseness=-1  
\end{abstract}

% The code is available at \href{https://github.com/sail-sg/Attention-Sink}{https://github.com/sail-sg/Attention-Sink}.

